\chapter{G6PD Enzyme Testing}
\section*{What Is This Test?}
Glucose-6-phosphate dehydrogenase (G6PD) testing measures the activity of the G6PD enzyme in red blood cells. G6PD catalyzes the first step of the pentose phosphate pathway, reducing NADP+ to NADPH, which is essential for protecting cells from oxidative stress. In red blood cells (which lack mitochondria and therefore lack the tricarboxylic acid cycle), the pentose phosphate pathway is the sole source of NADPH. NADPH maintains glutathione in its reduced state (GSH), which detoxifies hydrogen peroxide and organic peroxides via glutathione peroxidase, protecting hemoglobin and the RBC membrane from oxidative damage. G6PD deficiency is an X-linked recessive disorder (most common enzymopathy worldwide, affecting >400 million people, predominantly in malaria-endemic regions [Africa, Mediterranean, Middle East, Southeast Asia] because G6PD deficiency confers partial protection against falciparum malaria) \cite{luzzatto2020glucose,ruwende1998glucose}. In G6PD deficiency, exposure to oxidative stressors triggers acute hemolytic anemia: Heinz bodies (precipitated denatured hemoglobin) form, bite cells/blister cells appear on the peripheral smear, and hemoglobinemia/hemoglobinuria develop. Common triggers include: fava beans (favism --- the classic and most well-known trigger in the Mediterranean variant), infections (the most common trigger across all variants), and drugs (dapsone, primaquine, chloroquine, sulfonamides, nitrofurantoin, rasburicase, methylene blue, nalidixic acid, and phenazopyridine). The G6PD assay should NOT be performed during an acute hemolytic episode because reticulocytes (young RBCs) have higher G6PD activity than older RBCs; G6PD-deficient older RBCs have already been destroyed, leaving behind reticulocytes with near-normal G6PD activity, producing a false-negative result. The test should be repeated 2--3 months after the acute hemolytic episode has resolved.
\section*{Alternative Names}
G6PD; Glucose-6-Phosphate Dehydrogenase; G6PD Enzyme Activity; G6PD Deficiency Screen; G6PD Quantitative Assay
\section*{Principle}
G6PD activity is measured by a quantitative spectrophotometric assay. Whole blood (EDTA or heparin) is incubated with glucose-6-phosphate and NADP+. G6PD catalyzes the oxidation of glucose-6-phosphate to 6-phosphogluconate, and NADP+ is reduced to NADPH. The rate of NADPH production is measured by the increase in absorbance at 340 nm. G6PD activity is expressed in U/g hemoglobin or U/10$^{12}$ RBCs. Screening tests (fluorescent spot test) provide qualitative results: NADPH fluoresces under UV light; absence of fluorescence indicates G6PD deficiency.
\section*{History}
G6PD deficiency was discovered in 1956 during investigation of acute hemolytic reactions to the antimalarial primaquine in US soldiers in Korea, when Carson and colleagues showed that primaquine-sensitive individuals lacked glucose-6-phosphate dehydrogenase activity in their red blood cells \cite{carson1956enzymatic}. Ernest Beutler and coworkers introduced the glutathione stability test in 1957, enabling large-scale population screening, and the enzyme was soon shown to be X-linked, explaining its male predominance. Recognition that the deficiency underlies favism, the hemolysis after fava bean ingestion long known in Mediterranean countries, followed in the late 1950s. Over the 1960s--1970s, the classic G6PD variants (A-, Mediterranean) were defined by electrophoretic and biochemical characterization \cite{beutler1994g6pd}. Population genetic studies established that the geographic distribution of G6PD deficiency mirrors that of falciparum malaria, supporting the hypothesis that the deficiency confers partial protection against malaria, and molecular genotyping methods were introduced in the 1980s--1990s.
\section*{Why This Test Is Ordered}
\begin{itemize}
  \item Pre-treatment screening before prescribing oxidant drugs, especially primaquine, dapsone, rasburicase, sulfonamides, and nitrofurantoin
  \item Evaluation of acute hemolytic anemia after fava bean ingestion, infection, or drug exposure in patients from endemic regions
  \item Investigation of neonatal jaundice and hyperbilirubinemia in populations with a high prevalence of G6PD deficiency
  \item Evaluation of chronic non-spherocytic hemolytic anemia (CNSHA), which may represent severe G6PD variants
  \item Confirmation of G6PD deficiency after an episode of hemoglobinuria or unexplained acute anemia
  \item Family screening and genetic counseling, because the disorder is X-linked and carrier females may transmit it to sons
  \item Newborn screening in jurisdictions where G6PD deficiency is common \cite{who1989glucose}
\end{itemize}
\section*{Primary vs Secondary Test}
G6PD testing is a primary diagnostic test for the most common inherited red cell enzyme disorder, but it is ordered selectively when clinical suspicion is raised by hemolysis after an oxidative stressor or before the prescription of an oxidant drug. It is a screening and confirmatory test rather than a component of routine health panels.
\section*{How This Test Is Performed}
Whole blood is collected by standard venipuncture into an EDTA (purple-top) or heparinized tube. The quantitative spectrophotometric assay measures the rate of NADPH production at 340 nm when G6PD oxidizes glucose-6-phosphate in the presence of NADP+; the fluorescent spot test provides a qualitative screening result based on NADPH fluorescence under UV light \cite{beutler1968special}. Activity is expressed in U/g hemoglobin or U/10$^{12}$ red blood cells. Results are typically available within 1--2 days. The sample must be obtained during a steady state, not during an acute hemolytic episode.
\section*{What Preparation Is Needed}
None; fasting is not required. The test should not be collected during an acute hemolytic episode or within 2--3 months after resolution, because the reticulocytosis of recovery yields falsely normal activity. Recent red blood cell transfusion within the preceding 2--3 months also produces false-normal results and should be avoided if possible.
\section*{Contraindications and Precautions}
There are no absolute contraindications. Important cautions: (1) Do not test during acute hemolysis -- the reticulocyte response causes false-negative results; repeat testing should be deferred 2--3 months after resolution. (2) Recent transfusion contaminates the sample with donor red cells and masks deficiency. (3) Heterozygous females may have intermediate or normal activity because of variable X-inactivation and may require genotyping for diagnosis. (4) The fluorescent spot test is unreliable in individuals with high reticulocyte counts or severe anemia; the quantitative assay is preferred for borderline results.
\section*{Risks}
Minimal risk. Slight pain or bruising at the venipuncture site. Rarely: hematoma, infection, or vasovagal syncope.
\section*{What Happens After the Test}
The result is interpreted with the CBC, reticulocyte count, LDH, bilirubin, and peripheral smear to establish the hemolytic nature of the episode. A confirmed deficiency prompts patient education on avoidance of fava beans and oxidant drugs, prompt treatment of infections, and provision of a list of contraindicated medications; genotyping may be offered to define the variant and to identify carrier females. Neonates with G6PD deficiency and significant jaundice are managed with phototherapy and monitoring for kernicterus.
\section*{Understanding Results}
\subsection*{Normal G6PD Activity}
No G6PD deficiency. The patient is not at risk for oxidative hemolysis from primaquine, dapsone, or fava bean ingestion.
\subsection*{Mild to Moderate Deficiency}
G6PD activity 10--60\% of normal. African variant (G6PD A-): enzyme instability with normal catalytic activity in young RBCs, reduced in older RBCs. Hemolysis is usually mild and self-limited (affecting only older RBCs) after exposure to an oxidative stressor.
\subsection*{Severe Deficiency (<10\%)}
G6PD Mediterranean variant. Profound deficiency, severe acute hemolytic anemia after exposure to fava beans (favism), infections, or drugs \cite{cappellini2008glucose}. Can be life-threatening. Neonatal jaundice (severe hyperbilirubinemia and kernicterus risk).
\subsection*{False-Negative During Acute Hemolysis}
Reticulocytes have 5--10x higher G6PD activity than mature RBCs. During an acute hemolytic crisis, reticulocytosis yields falsely normal G6PD activity. The test should be repeated 2--3 months after resolution, or G6PD activity can be corrected for reticulocyte count.
\section*{Normal Results}
G6PD activity: 8.0--18.0 U/g hemoglobin (varies by laboratory). Quantitative: >60\% of normal = normal; 10--60\% = deficient; <10\% = severely deficient. Adjust for reticulocytosis during acute hemolysis.
\section*{Follow-up Tests}
Reticulocyte count (to assess for false-negative elevation), peripheral blood smear (blister cells, bite cells, Heinz bodies on supravital staining), G6PD genotyping (for specific G6PD variants), and repeat G6PD enzyme activity after 2--3 months.
\section*{References}
\bibliographystyle{unsrtnat}
\bibliography{references}
\clearpage
